\sechead{ОПЫТ РАБОТЫ}

\textbf{Яндекс / Международный поиск} \hfill \emph{06.2026 -- 08.2026}

\textit{Стажёр по машинному обучению}
\begin{cvitems}
    \item \textbf{Алерты и Alert Investigator:} самостоятельно реализовал и внедрил систему статистических алертов и LLM-агента для автоматического расследования отклонений продуктовых метрик. Агент анализирует проблемные диалоги, ищет паттерны, кластеризует кейсы и проверяет гипотезы. Используется командой для сужения круга ручных проверок.
    \item \textbf{Метрика успешности:} интегрировал LLM-метрику в A/B-тестирование, сделав её доступной разработчикам агентов для оценки продуктовых изменений. Адаптировал расчёт метрики для дополнительных пользовательских сценариев.
    \item \textbf{Развитие модели успешности:} провёл офлайн-эксперименты по улучшению качества оценки успешности. Проверял гипотезы об использовании дополнительного контекста и сравнивал варианты промпта.
    \item \textbf{Языковой классификатор:} разработал прототип и передал его разработчикам агентов. Провёл оценку качества и анализ ошибок. Собрал пайплайн подготовки данных для дистилляции.
    \item \textbf{Продуктовая аналитика:} разработал в DataLens дашборд для анализа связи времени ответа с поведением пользователей. Самостоятельно подготовил данные и настроил регулярный расчёт показателей. Предоставлял аналитические срезы нескольким продуктовым направлениям. Данные использовались командами агентов при обсуждении задержек ответа.
\end{cvitems}

\textbf{АО <<Площадь>>} \hfill \emph{09.2023 -- 01.2024}

\textit{Онлайн-стажировка -- сегментация лесных пожаров на видео с БПЛА.}\vspace{0.12cm}
	
\textbf{Сибирская генерирующая компания, Новосибирская ТЭЦ-2} \hfill \emph{03.2023 -- 07.2023}
	
\textit{Некоммерческий ML-проект -- ансамбль моделей для прогноза расхода топлива.}\vspace{0.12cm}
